\chapter{Data Quality and Lineage}\label{ch:pl:data-quality-and-lineage}

On 6 May 2010, between 2:40 and 3:00 in the afternoon, more than 20\,000 trades in more than 300 US securities were executed at
prices 60\% or more away from where they had been minutes before, some at a penny, some at \$100\,000. After the close the
exchanges and their self-regulator agreed to break all of them as clearly erroneous. A \omterm{def:pl:a-tick-store-on-open-formats:store}{tick store} that had captured those
trades faithfully, and never learned that they had been cancelled, fed them to every backtest, volatility estimate and
execution study that touched that afternoon, for as long as it kept the data. Nothing in the store was corrupt; the data were
exactly what the feed had sent. What was missing was a layer that checks data against what it should look like, marks what
fails instead of deleting it, applies the corrections that arrive later as new versions, and knows which results were built
from what. This chapter builds that layer (\texttt{firm.dataqual}) and measures it on a simulated day with planted defects and,
just as important, genuine market events that a careless rule mistakes for defects.

\section{Checks: completeness, validity, consistency, timeliness}

\begin{definition}[Data-quality rule, data contract]\label{def:pl:data-quality-and-lineage:rule}
A \emph{data-quality rule}\index{data-quality rule} is an automated check of a dataset against an expectation --- that it is
complete, that its values are valid, that it is consistent with itself and with other sources, that it is timely --- run every
time the data are produced, with a severity and an owner who acts on its failures. A \emph{data contract}\index{data contract}
is the agreement between a dataset's producer and its consumers that states its schema, its rules and its delivery times, so
that a change on either side is a change to the contract rather than a surprise.
\end{definition}

Four families of checks cover most market-data defects. Completeness: no messages missing (the feed's sequence numbers have no
gap, chapter~2), no instrument silent when it should be quoting. Validity: every value in its domain (positive prices and sizes,
known instruments, timestamps inside the session). Consistency: the data agree with themselves (no bid at or above the ask, trade
prices near the quotes) and with other sources (another feed, a vendor's copy, the official close). Timeliness: the data arrived
when the contract says. \Cref{lst:pl:data-quality-and-lineage:rules} is the suite of the chapter, written as data: each rule has
a kind, a severity, an owner and parameters, and \texttt{firm.dataqual} runs any suite over any partition.

\omcode{code/platforms/07-data-quality-and-lineage/python/pl_dataqual.py}{88}{110}{The chapter's rule suite as data: sequence gaps,
bounds, crossed books, staleness, price spikes, broken trades and a count against a vendor's copy, each with a severity and an
owner.}\label{lst:pl:data-quality-and-lineage:rules}

The day is the second of the chapter-4 week: two instruments, 47\,472 events, 9\,258 quotes and 1\,190 trades. Two genuine events
are part of it: SIM2's trading is halted for 60 seconds at 40\% of the session, and SIM2's price level jumps by 3\% at 60\% of it,
on news, and stays there. Into copies of the day, twenty times with twenty seeds, the chapter plants six kinds of defect: a gap of
200 events, three crossed quotes, eight trade prices off by a factor of ten (a decimal slip, up or down), a 90-second silence of
SIM1's quotes, four trades broken by the venue after the close, and, separately, a vendor's later correction of an official close.

\section{Anomaly flags, not deletions}

\begin{definition}[Anomaly flag, quarantine]\label{def:pl:data-quality-and-lineage:flag}
An \emph{anomaly flag}\index{anomaly flag} is a record attached to a row or a partition that names the rule it failed, the
severity and the detail, leaving the row itself unchanged; consumers decide whether to exclude flagged rows. A
\emph{quarantine}\index{quarantine} is the state of a partition that failed an error-severity rule: it is withheld from readers
until a person releases it, with a reason, or it is rebuilt.
\end{definition}

Deleting a bad row is tempting and wrong. The row may be right (a genuine 3\% jump looks like a spike to a naive rule); another
consumer may need it (a surveillance system wants the erroneous trades, chapter~30); and a deletion cannot be audited or undone.
A flag keeps all three options open: the backtest excludes error-severity rows, the surveillance system reads them, and the rule
can be improved and rerun.

\begin{omfigure}
\begin{tikzpicture}[font=\scriptsize,>=Stealth,box/.style={draw,rounded corners=2pt,fill=omDef!10,minimum height=0.7cm,inner sep=3pt,align=center}]
\node[box] (p) at (0,0) {partition\\(date, table)};
\node[box,fill=omMeth!15] (r) at (2.9,0) {rule suite\\(data)};
\node[box,fill=omThm!12] (f) at (5.7,0) {flags\\(row, rule, severity)};
\node[box] (pub) at (9.2,0.8) {published:\\consumers filter};
\node[box,fill=omThm!20] (qz) at (9.2,-0.8) {quarantined:\\withheld};
\node[box,fill=gray!10] (rel) at (12.1,-0.8) {released\\by a person};
\node[box,fill=omProp!15] (cor) at (2.9,-1.8) {vendor correction:\\new version};
\draw[->] (p) -- (r); \draw[->] (r) -- (f);
\draw[->] (f) -- node[above,sloped]{warnings} (pub.west);
\draw[->] (f) -- node[below,sloped]{errors} (qz.west);
\draw[->] (qz) -- (rel);
\draw[->] (cor) -| (p);
\end{tikzpicture}

\omcaption{The data-quality layer: every partition runs the rule suite; flags are recorded beside the rows, never in place of them;
a partition with an error-severity flag is quarantined until released; a vendor's correction arrives as a new version of the data,
and the rules run again.}\label{fig:pl:data-quality-and-lineage:pipeline}
\end{omfigure}

\begin{omfigure}
\begin{tikzpicture}
\begin{axis}[width=12cm,height=5.6cm,ymode=log,xlabel={SIM2 trade number, in sequence order},ylabel={trade price (log scale)},
  tick label style={font=\small},label style={font=\small},grid=major,grid style={gray!20},ymin=20,ymax=5000,
  ytick={30,100,300,1000,3000},yticklabels={30,100,300,1\,000,3\,000},
  legend columns=2,legend style={font=\footnotesize,at={(0.5,-0.3)},anchor=north,draw=none,/tikz/every even column/.append style={column sep=8pt}},
  table/col sep=comma]
\addplot[only marks,mark=*,mark size=0.8pt,omDef] table[x=n,y=price]{figdata/platforms/07-data-quality-and-lineage/sim2_prices.csv};
\addplot[only marks,mark=o,mark size=3pt,omThm,thick] table[x=n,y=price,restrict expr to domain={\thisrow{flagged}}{1:1}]{figdata/platforms/07-data-quality-and-lineage/sim2_prices.csv};
\legend{trade,flagged as a spike (error)}
\end{axis}
\end{tikzpicture}

\omcaption{SIM2's trades on the defective copy of seed~0, on a log scale: two decimal slips (one ten times too high, one ten times too
low) stand out and are flagged; the genuine 3\% jump late in the day is invisible at this scale and is not flagged by the aware rule.
Data: \texttt{fig\_dataqual.py}.}\label{fig:pl:data-quality-and-lineage:prices}
\end{omfigure}

The first version of the suite flagged the planted defects and a good deal more. The spike rule compared each trade price with the
median of the previous 50 trades of its symbol, in robust units: the median absolute deviation of the previous deviations, with a
floor of 0.1\% so that a one-tick move is never a spike. On the day of the jump it flagged every trade for the next few dozen, until the backward median caught up with the
new level (\cref{fig:pl:data-quality-and-lineage:jump}). A spike and a level shift look identical from behind; from ahead they do
not, because a spike comes back. The data-quality run happens after the close, so it may look ahead:

\begin{method}[Telling a spike from a move]\label{met:pl:data-quality-and-lineage:spike}
Flag a value when it deviates, by more than $k$ robust units, both from the median of the values before it and from the median of
the values after it. A decimal slip deviates from both; a genuine move of the level deviates only from the values before. At the
start of a series, where there is no history, compare with the values after only.
\end{method}

\omcode{code/firm/dataqual/firm_dataqual.py}{78}{99}{The spike rule: deviation from the backward median in robust units, floored,
and, with \texttt{confirm}, from the forward median too.}\label{lst:pl:data-quality-and-lineage:spike}

\begin{omfigure}
\begin{tikzpicture}
\begin{axis}[width=12cm,height=5.8cm,xlabel={SIM2 trade number},ylabel={price},tick label style={font=\small},
  label style={font=\small},grid=major,grid style={gray!20},enlarge y limits=0.12,
  legend columns=2,legend style={font=\footnotesize,at={(0.5,-0.3)},anchor=north,draw=none,/tikz/every even column/.append style={column sep=8pt}},
  table/col sep=comma]
\addplot[only marks,mark=*,mark size=1.2pt,omDef] table[x=n,y=price]{figdata/platforms/07-data-quality-and-lineage/jump.csv};
\addplot[thick,omMeth] table[x=n,y=behind]{figdata/platforms/07-data-quality-and-lineage/jump.csv};
\addplot[thick,omMeth,dashed] table[x=n,y=ahead]{figdata/platforms/07-data-quality-and-lineage/jump.csv};
\addplot[only marks,mark=o,mark size=2.6pt,omThm,thick] table[x=n,y=price,restrict expr to domain={\thisrow{naive_only}}{1:1}]{figdata/platforms/07-data-quality-and-lineage/jump.csv};
\legend{trade price,median of the previous 50,median of the next 50,flagged by the naive rule only}
\end{axis}
\end{tikzpicture}

\omcaption{SIM2's trades around its genuine 3\% jump. Measured from the backward median, every trade for the next few dozen looks
like a spike, and the naive rule flags 24 of them (circled); measured from the forward median too, none does, since the new level
persists. Data: \texttt{fig\_dataqual.py} (seed 0).}\label{fig:pl:data-quality-and-lineage:jump}
\end{omfigure}

The staleness rule failed the same way on the halt: 60 seconds without a quote from SIM2 is a defect in the feed only if the market
was open. The rule must know the market's state --- the venue's trading-action messages (chapter~2's events include them) --- and
not count time spent halted (\cref{lst:pl:data-quality-and-lineage:stale}).

\omcode{code/firm/dataqual/firm_dataqual.py}{102}{117}{The staleness rule, which does not count the time the market itself spent
silent (a trading halt, a scheduled pause) in a gap.}\label{lst:pl:data-quality-and-lineage:stale}

\Cref{tab:pl:data-quality-and-lineage:detection} is the result over twenty seeds. Every planted defect is found by both versions of
the suite; the difference is entirely in the false flags: 495 trades wrongly flagged as spikes and 20 halts wrongly flagged as
stale feeds by the naive suite, none by the aware one.

\begin{table}[ht]
\centering\small
\begin{tabular}{lrrrr}
\toprule
rule & planted & found & false flags, naive & false flags, aware \\
\midrule
sequence gap & 20 & 20 & 0 & 0 \\
crossed book & 60 & 60 & 0 & 0 \\
price spike & 160 & 160 & 495 & 0 \\
stale quotes & 20 & 20 & 20 & 0 \\
broken trades & 80 & 80 & 0 & 0 \\
\bottomrule
\end{tabular}
\caption{Planted defects, defects found and rows wrongly flagged by each rule over twenty seeds of planting, with the naive suite
(backward-looking spike rule, staleness blind to halts) and the aware one. Data: \texttt{pl\_dataqual.detection}.}\label{tab:pl:data-quality-and-lineage:detection}
\end{table}

A planted defect that every rule finds is not a test of the rules; the confounders are. A surveillance detector of One Quant
Book~9 caught all its planted spoofers until legitimate look-alikes were added (its chapter~29), and the same lesson holds here:
plant the genuine events the real market has, and measure the false flags.

\section{Cross-source checks}

\begin{definition}[Cross-source check]\label{def:pl:data-quality-and-lineage:cross}
A \emph{cross-source check}\index{cross-source check} compares a dataset with an independent source of the same facts --- the other
feed line, a second vendor, the venue's end-of-day file, the official close --- on an aggregate that both can compute (a count, a
volume, a last price per interval), and flags the intervals where they differ by more than a tolerance.
\end{definition}

A single source can be complete by its own sequence numbers and still wrong: a normaliser that dropped a message type, a
compaction that lost a file. The chapter's check counts trades per minute in the store and in a vendor's copy of the day and flags
minutes that differ by more than 2\%: over the twenty seeds it raised 16 flags, the minutes in which the planted gap of 200 events
happened to contain trades. It does not say which trades are missing, only where to look; that is the right division of labour
between a cheap check that runs everywhere and an expensive investigation that runs where it points.

\section{Vendor corrections and versions}

\begin{definition}[Vendor correction]\label{def:pl:data-quality-and-lineage:correction}
A \emph{vendor correction}\index{vendor correction} is a later delivery that changes a value already delivered (a close, a trade
condition, a corporate action): it is stored as a new version of the same fact, with its own knowledge time, and never overwrites
the earlier version.
\end{definition}

The chapter's vendor delivers SIM2's official close at 16:30 on the day, and a corrected close at 07:10 the next morning. Both are
versions of one fact in Book~7's bitemporal store: asked at 23:00, the close is the first; asked at 08:00, the second. The broken
trades of the hook are the same pattern: a list of cancellations, delivered after the close, applied as flags (\texttt{busted}) on
the trades they name. What a correction changes is not only the fact but everything computed from it, and that is lineage.

\section{Lineage and provenance}

\begin{definition}[Data lineage, provenance record]\label{def:pl:data-quality-and-lineage:lineage}
The \emph{data lineage}\index{data lineage} of a result is the graph of dataset versions and transformations it was derived from, back
to its raw inputs. A \emph{provenance record}\index{provenance record} is the stored description of how one dataset version was
produced: the versions of its inputs (by content hash), the code and its version, and when and by whom it ran.
\end{definition}

\texttt{firm.dataqual}'s lineage records each dataset version as a node, named by the content hash of its data (Book~7's
\texttt{firm.workflow.content\_hash}), with an edge from each input version it was built from
(\cref{lst:pl:data-quality-and-lineage:lineage}). A correction produces a new version of the input; every node downstream of the
old version is stale. In the chapter's example (\cref{fig:pl:data-quality-and-lineage:lineage}), the corrected close makes the
marks, the P\&L and the research features stale, and leaves the minute bars, built from the trades alone, untouched.

\omcode{code/firm/dataqual/firm_dataqual.py}{189}{203}{Lineage: a dataset version is its name and content hash, with its inputs'
versions; everything downstream of an old version is stale once it is corrected.}\label{lst:pl:data-quality-and-lineage:lineage}

\begin{omfigure}
\begin{tikzpicture}[font=\scriptsize,>=Stealth,box/.style={draw,rounded corners=2pt,fill=omDef!10,minimum height=0.6cm,inner sep=3pt,align=center}]
\node[box] (tr) at (0,0) {trades\\2026-09-22};
\node[box,fill=omMeth!25] (cl) at (3,0.9) {closes\\(corrected)};
\node[box] (ba) at (3,-0.9) {minute bars};
\node[box,fill=omThm!12] (mk) at (6,1.5) {marks};
\node[box,fill=omThm!12] (pn) at (9,1.5) {P\&L};
\node[box,fill=omThm!12] (ft) at (6,-0.3) {features};
\draw[->] (tr) -- (cl); \draw[->] (tr) -- (ba);
\draw[->] (cl) -- (mk); \draw[->] (mk) -- (pn);
\draw[->] (cl) -- (ft); \draw[->] (ba) -- (ft);
\node[omThm,anchor=west] at (7.2,-0.3) {stale: rebuild};
\node[anchor=north] at (3,-1.25) {not affected};
\end{tikzpicture}

\omcaption{The lineage of the day's results. The vendor's correction of the close makes every version built from the old close stale
(red): the marks, the P\&L and the research features; the minute bars, built from the trades alone, are not affected. Data:
\texttt{pl\_dataqual.correction\_and\_lineage}.}\label{fig:pl:data-quality-and-lineage:lineage}
\end{omfigure}

\begin{dated}{2026-09}{A standard for provenance}\label{dat:pl:data-quality-and-lineage:prov}
The World Wide Web Consortium's PROV data model (a W3C Recommendation of 30 April 2013, consulted September 2026) defines provenance
as a record that describes the people, institutions, entities and activities involved in producing, influencing or delivering a
piece of data; the lineage graph of this chapter is a small instance of its entities (dataset versions) and activities (the jobs
that derive them).
\end{dated}

What quality is worth shows in the numbers of the day (seed~0). With the planted defects kept, SIM1's VWAP is 2.69\% too high and
SIM2's 1.08\%, and SIM2's realised variance of trade-price changes is 24\,000 times its true value: eight decimal slips outweigh a
day of genuine trading. With the error-severity rows excluded, both VWAPs are within 0.03\% of the truth (the residue is the trades
lost in the gap, which no flag can restore) and the realised variance equals the true one.

\section{Tutorial: plant, flag, correct}\label{tut:pl:data-quality-and-lineage:tut}

\begin{tutorial}
\textbf{Goal.} Plant defects and confounders in a simulated day, run the rule suite twice (naive and aware), measure detection and
false flags, and trace a correction through the lineage. \textbf{End state:}
\cref{tab:pl:data-quality-and-lineage:detection}, \cref{fig:pl:data-quality-and-lineage:jump} and the VWAP and variance errors of the
last section.
\begin{tutsteps}
\item \textbf{The day}: \texttt{pl\_dataqual.base()} with its halt and its jump.
\item \textbf{Plant}: \texttt{plant(seed)} returns the defective tables and the truth of which rows are defects.
\item \textbf{Run}: \texttt{firm\_dataqual.run(rules(\dots), tables)} (\cref{lst:pl:data-quality-and-lineage:rules}); then
\texttt{score} and \texttt{detection} over twenty seeds, naive and aware.
\item \textbf{Effect}: \texttt{effect(0)}: VWAP and realised variance kept, flagged out and true.
\item \textbf{Correct}: \texttt{correction\_and\_lineage()}: two versions of the close, and the stale results.
\end{tutsteps}
\textbf{What to change next.} Lower the spike threshold $k$ from 8 to 4 and count the false flags on the natural day; plant a 2\%
spike instead of a decimal slip and find the smallest one the rule still catches.
\end{tutorial}

\section{Build: the data-quality layer}\label{bld:pl:data-quality-and-lineage:dataqual}

\begin{build}
\bfield{Purpose} The checks every partition of the \omterm{def:pl:a-tick-store-on-open-formats:store}{tick store} passes before it is published, the flags consumers filter on, the
\omterm{def:pl:data-quality-and-lineage:flag}{quarantine} that withholds bad partitions, and the lineage that says what a correction invalidates; it runs as jobs of chapter~1's
map and chapter~13's scheduler.
\bfield{Interface} \texttt{Rule(name, kind, severity, owner, params)}, \texttt{Flag(table, row, rule, severity, detail)};
\texttt{run(rules\_by\_table, tables)}, \texttt{clean(df, flags, table, severities)}; \texttt{Quarantine} with \texttt{apply},
\texttt{release}, \texttt{readable}, \texttt{log}; \texttt{Lineage} with \texttt{add}, \texttt{downstream}, \texttt{upstream},
\texttt{stale\_after}; kinds \texttt{sequence}, \texttt{bounds}, \texttt{crossed}, \texttt{spike}, \texttt{stale},
\texttt{cross\_source}, \texttt{busted}.
\bfield{Rules} Rules are data with an owner; flags never modify rows; an error flag \omterm{def:pl:data-quality-and-lineage:flag}{quarantines} its partition; a release names a
person and a reason; corrections are versions; dataset versions are named by content hash.
\bfield{Acceptance tests} \texttt{code/firm/dataqual/tests/}: each kind finds its planted defect and nothing else; the confirmed spike
rule spares a level shift and the staleness rule a halt; the cross-source count by bucket; \omterm{def:pl:data-quality-and-lineage:flag}{quarantine} and release; lineage names the
stale results of a correction.
\bfield{Stretch} Rules generated from a \omterm{def:pl:data-quality-and-lineage:rule}{data contract}; a quality score per partition published with it; flags stored as a table of
the \omterm{def:pl:a-tick-store-on-open-formats:store}{tick store}, queryable like the data.
\end{build}

\begin{omsources}
\item U.S. Commodity Futures Trading Commission and Securities and Exchange Commission, \emph{Findings regarding the market events
of May~6, 2010}, 30 September 2010.
\item W3C, \emph{PROV-DM: The PROV Data Model}, W3C Recommendation, 30 April 2013.
\end{omsources}

\section{Exercises}

\begin{exercise}[$\star$]\label{exo:pl:data-quality-and-lineage:1}
A feed's sequence numbers jump from 91\,686 to 91\,691. How many messages are missing, and what does the capture of chapter~2 do
about them?
\end{exercise}

\begin{exercise}[$\star$]\label{exo:pl:data-quality-and-lineage:2}
Why does the spike rule have a floor on its scale? What would happen without one on a large-tick instrument whose trades mostly
print at the same price?
\end{exercise}

\begin{exercise}[$\star$]\label{exo:pl:data-quality-and-lineage:3}
A trade prints at 3\,000\,000 when the previous trades were at 300\,000 and the next ones are too. What kind of defect is it likely
to be, and why does the forward median confirm it?
\end{exercise}

\begin{exercise}[$\star\star$]\label{exo:pl:data-quality-and-lineage:4}
Why should a rule flag and not delete? Give one consumer that needs the flagged rows.
\end{exercise}

\begin{exercise}[$\star\star$]\label{exo:pl:data-quality-and-lineage:5}
With the defects kept, SIM2's VWAP on seed~0 is 1.08\% too high; after excluding error rows it is 0.02\% too low. Where does each
error come from?
\end{exercise}

\begin{exercise}[$\star\star$]\label{exo:pl:data-quality-and-lineage:6}
The official close is corrected at 07:10 the next morning. Which of the day's results must be rebuilt, which need not, and how does
the lineage know?
\end{exercise}

\begin{exercise}[$\star\star\star$]\label{exo:pl:data-quality-and-lineage:7}
\emph{Coding.} Run \texttt{score} for seed~0 with the naive and the aware suite. How many false flags does each rule raise, and
which genuine event causes each?
\end{exercise}

\begin{exercise}[$\star\star\star$]\label{exo:pl:data-quality-and-lineage:8}
\emph{Find the flaw.} ``Our data-quality job deletes any trade more than 5\% from the previous trade and any gap in quotes longer than
30 seconds is filled by repeating the last quote. The cleaned data go straight to the \omterm{def:pl:a-tick-store-on-open-formats:store}{tick store}.''
\end{exercise}

\section{Problem: The Afternoon That Was Cancelled}

\begin{problem}[{Weekend problem --- a day, its defects and its genuine surprises}]\label{pb:pl:data-quality-and-lineage:1}
The chapter's simulated day (a halt and a 3\% jump of SIM2), its six planted defects and the two rule suites.

\medskip\noindent\textbf{Part I --- The checks.}
\begin{enumerate}
\item Name the four families of checks and give a rule of each from the chapter's suite.
\item What is a \omterm{def:pl:data-quality-and-lineage:rule}{data contract}, and what does it add to the rules?
\item How many events, quotes and trades does the day have?
\item Which of the planted defects can no rule on the day's own data detect, and what detects it?
\item Why is the vendor's close correction not a data-quality flag?
\end{enumerate}

\medskip\noindent\textbf{Part II --- False flags.}
\begin{enumerate}[resume]
\item Why does the naive spike rule flag trades after the 3\% jump, and how many on seed~0?
\item How many false spike flags over twenty seeds, naive and aware?
\item Why does the naive staleness rule flag the halt, and how does the aware rule avoid it?
\item What would a pipeline that deleted flagged rows have done to the day of the jump?
\item Why must a test of \omterm{def:pl:data-quality-and-lineage:rule}{data-quality rules} include genuine events, not only planted defects?
\end{enumerate}

\medskip\noindent\textbf{Part III --- What quality is worth.}
\begin{enumerate}[resume]
\item How far are the two VWAPs from the truth with the defects kept, and with error rows excluded?
\item By what factor is SIM2's realised variance inflated with the defects kept?
\item What does the \omterm{def:pl:data-quality-and-lineage:cross}{cross-source check} find over twenty seeds, and what does it not say?
\item How does a \omterm{def:pl:data-quality-and-lineage:flag}{quarantine} differ from a flag, and who releases it?
\item Which results does the corrected close make stale?
\end{enumerate}

\medskip\noindent\textbf{Part IV --- The verdict.}
\begin{enumerate}[resume]
\item State the \emph{named result}: the detection and false-flag counts of the two suites over twenty seeds, and the VWAP and variance
errors with and without flags.
\item Which single change to the naive suite removed the most false flags?
\item What would the store have held for 6 May 2010 afternoon, with and without the broken-trade rule?
\item What must the lineage record so that a correction names every stale result?
\item In one sentence: what distinguishes a data-quality layer from a cleaning script?
\end{enumerate}
\end{problem}

\section{Interview questions}

\begin{interviewq}[$\star$ \iqroles{developer, researcher}]\label{iq:pl:data-quality-and-lineage:1}
What checks would you run on a day of \omterm{def:pl:capturing-and-storing-tick-data:tick}{tick data} before researchers use it?
\end{interviewq}

\begin{interviewq}[$\star\star$ \iqroles{developer, researcher}]\label{iq:pl:data-quality-and-lineage:2}
How do you tell a bad print from a genuine price move?
\end{interviewq}

\begin{interviewq}[$\star\star$ \iqroles{developer}]\label{iq:pl:data-quality-and-lineage:3}
Why flag bad data rather than delete it?
\end{interviewq}

\begin{interviewq}[$\star\star$ \iqroles{developer, risk}]\label{iq:pl:data-quality-and-lineage:4}
A vendor corrects yesterday's closing price. What has to happen in the firm's systems?
\end{interviewq}

\begin{interviewq}[$\star\star\star$ \iqroles{developer}]\label{iq:pl:data-quality-and-lineage:5}
How would you test a data-quality system, and how would you know it is not flagging too much?
\end{interviewq}

\begin{interviewq}[$\star\star\star$ \iqroles{developer}]\label{iq:pl:data-quality-and-lineage:6}
Design lineage for a research platform so that any published result can be traced to its inputs and recomputed after a correction.
\end{interviewq}
